%% ======================================================================
%% 10_conclusions.tex -- 9. Conclusiones
%% ======================================================================
\section{Conclusiones}
\label{sec:conclusions}
AeroTray muestra que tres lazos de equilibrio independientes pueden apilarse sobre una base móvil con hardware de bajo costo, y que la propagación de la perturbación de la base a través de ellos puede tanto medirse con registros sincronizados como predecirse con un modelo de primeros principios. El modelo, verificado simbólicamente y parametrizado con el ensamble CAD, las hojas de datos y los registros, lleva a seis conclusiones que los experimentos confirman.

(i) El módulo aéreo es un péndulo invertido sobre un pivote acelerado cuyo centro de masa, identificado con telemetría de vuelo estacionario, está $0{,}14\pm\SI{0.06}{\metre}$ sobre la rótula y \SI{9.5}{\milli\metre} hacia M8. (ii) El lazo de actitud funciona solo porque la rigidez proporcional supera a la gravitatoria por un margen pequeño, que proporcionó la ganancia de empuje subestimada; ese margen explica el modo lento de actitud observado en vuelo estacionario. (iii) Las divergencias lentas de los primeros vuelos eran un problema de autoridad de par, con una región de atracción de solo 3--4$^\circ$, y se resolvieron ampliando el límite de la velocidad angular pedida, no reajustando ganancias. (iv) El lazo de la bola está limitado por su retardo y por un cero en el semiplano derecho que introduce la altura de la bola sobre el pivote del plato, que juntos explican el techo de la ganancia derivativa; la ganancia efectiva de la planta es alrededor del 70\% de la del CAD por el juego del mecanismo. (v) La rótula pasiva elimina la pendiente de la dinámica de la bola, los giros son benignos, y el error residual de la bola con los rotores girando lo domina una carga aerodinámica y vibratoria directa y lentamente variable, no el movimiento de actitud del octocóptero.

(vi) El dron se justifica en la rampa por dos razones cuantificadas. Al nivelarse, desacopla la pendiente del plato y deja libre entre el 84 y el 95\% del recorrido de los servos, que con la bandeja fija se agotaría en una pendiente de \SI{11}{\degree}. Al sostener parte de su peso, actúa como un lastre ajustable que baja el centro de masa: con la mitad del peso del módulo en los rotores, la pendiente de vuelco al subir pasa de \SI{21.5}{\degree} (medida con el robot completo de \SI{10.4}{\kilo\gram}) a \SI{27.7}{\degree} (prevista), al bajar de \SI{30.5}{\degree} a \SI{38.2}{\degree} (prevista), y el margen de aceleración en \SI{20}{\degree} se multiplica por 3,5. El dron no debe flotar en la rampa, sino trabajar en una ventana del 35--55\% de su peso, a un costo de 2--5\% de batería por minuto. El margen de empuje para flotar, en cambio, depende del cuadrado de la tensión de las LiPo: con 12,4\,V el módulo despega a 1725--\SI{1775}{\micro\second}, y con 11,1\,V necesitaría 1850--\SI{2000}{\micro\second}, lo que explica que dejara de flotar y limita el vuelo útil a 3--\SI{7}{\minute} por par de baterías.

Experimentalmente, el firmware final mantiene el módulo de \SI{3.6}{\kilo\gram} en la guía con 0,55--\SI{1.7}{\degree} de error RMS de actitud, y el controlador de la bola la mantiene a una mediana de 19--\SI{30}{\milli\metre} del centro con la base quieta, en llano y en una rampa de \SI{10}{\degree}. En los ensayos finales sobre el robot completo (23/11) se registraron 13 ensayos de equilibrio de 23--\SI{122}{\second} (\SI{816}{\second} totales; mediana de 29--\SI{44}{\milli\metre}), con los rotores encendidos en 10 de ellos y el dron nivelado a 0,71--\SI{1.0}{\degree} RMS; en cuatro ensayos completos la base cruzó la rampa de \SI{10}{\degree} con los rotores encendidos y la bola en equilibrio durante 18--\SI{30}{\second} (mediana de 30--\SI{35}{\milli\metre}, dron a 0,85--\SI{1.15}{\degree} RMS sin ningún corte, con la bola saliendo del plato en el tramo final), y el módulo también despegó y aterrizó en la guía con la bandeja encima (la bola salió al despegar). La velocidad de la base medida ese día coincide con el modelo dentro de un 8--10\%. Dos correcciones concretas salen de esos ensayos: volver a nivelar el neutro del eje $x$ del plato, que concentra el error, y suavizar el despegue en la guía, cuyo salto por fricción es lo que hizo caer la bola al despegar. El modelo también señala los siguientes pasos con mayor beneficio esperado: medir la tensión de las baterías y compensar el colectivo, ensayar el modo rampa, bajar el centro de masa del módulo por debajo de la rótula, acortar el retardo del lazo de la bola, proteger el plato del flujo de los rotores y suavizar las referencias de la base. Todos los archivos de diseño, el código, los scripts de análisis, los videos y los datos crudos se publican para que estos experimentos puedan reproducirse y ampliarse, por ejemplo con la prealimentación de inclinación ya implementada o con controladores entrenados sobre el modelo publicado.
